\chapter{Introducción}
\label{cap:introduccion}

Este trabajo nace de una observación sencilla: los hogares inteligentes de hoy no son tan inteligentes como parecen. La mayoría de los dispositivos domóticos dependen de servidores externos para funcionar, de modo que cuando falla la conexión a Internet, o cuando el fabricante decide dejar de dar soporte, el usuario se queda sin poder encender una luz o regular el termostato.

El problema no es menor. Cada vez que un sensor de temperatura envía datos a la nube, o que un asistente de voz procesa una orden en servidores remotos, la privacidad del usuario queda expuesta. Los patrones de comportamiento doméstico (a qué hora se levanta, qué temperatura prefiere para dormir, qué días está en casa) viajan a servidores de terceros, a menudo sin que el usuario sea consciente ni haya dado un consentimiento informado.

El proyecto que se presenta en esta memoria, \textbf{Luxe Core AI}, aborda estos problemas desde un enfoque distinto: en lugar de delegar la inteligencia del hogar a la nube, la ejecuta íntegramente en hardware local. El sistema no necesita Internet para funcionar. Los modelos de lenguaje que toman decisiones se ejecutan en un servidor doméstico, y los datos de los sensores nunca abandonan la red local.

Esta decisión de diseño tiene consecuencias prácticas. La latencia de respuesta se reduce de segundos a milisegundos para la mayoría de los comandos. El sistema sigue funcionando aunque el router pierda la conexión con el exterior. Y, sobre todo, los datos del hogar pertenecen exclusivamente a quien vive en él.

Una de las características más singulares del ecosistema es que integra dispositivos que no fueron diseñados para interoperar. El proyecto incluye un enchufe inteligente Tuya controlado sin pasar por la nube del fabricante (que a su vez alimenta un extractor de humedad para pecera, un dispositivo originalmente ``tonto'' que solo necesita corriente y que, al conectarlo al enchufe, adquiere control remoto), un sensor de temperatura Xiaomi leído mediante Bluetooth Low Energy sin necesidad de su aplicación oficial, y un ventilador de techo de 35~\euro{} adquirido en AliExpress cuyo protocolo de comunicación fue descubierto mediante ingeniería inversa. Todos estos dispositivos (que individualmente son ``tontos'') adquieren capacidad de razonamiento al ser orquestados por la capa de inteligencia artificial. El resultado es un ecosistema donde la interoperabilidad no depende de estándares cerrados ni de la buena voluntad de los fabricantes, sino de una capa de abstracción inteligente capaz de entender y coordinar hardware heterogéneo.

El desarrollo ha seguido dos fases claramente diferenciadas. Durante la primera, se utilizó la infraestructura de la Universidad de Córdoba (servidores con GPU RTX accesibles por VPN) para prototipar la arquitectura de agentes y validar que el enfoque era viable. En la segunda fase, todo el sistema migró a una estación de trabajo doméstica, eliminando cualquier dependencia externa. Esta evolución no fue un retroceso, sino la materialización del objetivo central del proyecto: demostrar que la domótica inteligente puede funcionar sin depender de nadie más que del propio usuario.

\section{Motivación Personal}

Más allá de las consideraciones técnicas, este proyecto responde a una motivación personal: aprender a trabajar con tecnologías de \textit{Edge AI} en un entorno real, con hardware real, y comprobar por mí mismo hasta dónde pueden llegar los modelos de lenguaje pequeños ejecutados en local.

Cuando empecé el TFG, los modelos de lenguaje grandes (más de 30~B parámetros) copaban todos los titulares, pero requerían GPUs de varios miles de euros que estaban fuera de mi alcance. En paralelo, modelos mucho más ligeros (como \texttt{qwen2.5-coder:1.5b} o \texttt{bge-m3}) estaban demostrando que se podía hacer mucho con muy poco: clasificación de comandos, análisis semántico, generación de texto coherente, todo en menos de 500~MB de RAM. Quería comprobar si esos modelos eran suficientemente buenos para un caso de uso real, con las limitaciones de un servidor doméstico.

Los resultados han superado mis expectativas. El clasificador 1.5B responde en 300~ms, el razonador 7B conversa con naturalidad, y el sistema completo funciona 24/7 sin intervención. Esta experiencia me ha convencido de que los modelos pequeños y locales no son una solución temporal o de compromiso: son el futuro de la inteligencia artificial aplicada. Seguirán mejorando en precisión, velocidad y capacidad, y acabarán convertidos en una tecnología tan ubicua y asentada como lo es hoy el procesamiento en el \textit{edge} para tareas de visión artificial.

A nivel de documentación, esta memoria se ha redactado con \LaTeX{} y se gestiona con el mismo control de versiones (Git) que el código del proyecto. Esto permite mantener la trazabilidad entre las decisiones de diseño, el código que las implementa y la documentación que las describe.

Luxe Core AI aspira a ser, en definitiva, un ejemplo de que la inteligencia artificial en el hogar no tiene por qué implicar una pérdida de privacidad ni una dependencia irreversible de infraestructuras externas. Es una contribución (todavía modesta, pero funcional) a la idea de que la tecnología doméstica puede ser potente y soberana al mismo tiempo.

Cabe señalar que la estructura de esta memoria no sigue estrictamente la guía de desarrollo oficial del GII~\cite{uco_tfg_guide}, pues se ha optado por una organización temática que facilita la lectura del proceso completo: definición del problema, antecedentes, metodología, recursos, requisitos, especificación del sistema, validación y conclusiones. Esta decisión está alineada con la naturaleza del proyecto, que integra múltiples tecnologías y requiere un hilo narrativo continuo para su comprensión.

